%%%%%%%%%%%%%%%%%%%%%%%%%%%%%%%%%%%%%%%%%%%%%%%%%%%%%%%%%%%%
% 5.11 HARMONIC ANALYSIS
% The Composition of Advanced Mathematics
%%%%%%%%%%%%%%%%%%%%%%%%%%%%%%%%%%%%%%%%%%%%%%%%%%%%%%%%%%%%

\section{Harmonic Analysis}
\label{sec:harmonic-analysis}

\prerequisite{
Real analysis, measure theory, Lebesgue integration, \(L^p\) spaces,
complex analysis, functional analysis, Fourier series, Fourier transforms,
convolution, and basic probability.
}

\learningobjectives{
By the end of this section, the reader should be able to:
\begin{itemize}
    \item explain how harmonic analysis extends Fourier analysis;
    \item work with translation, dilation, modulation, and convolution;
    \item use approximate identities to smooth and recover functions;
    \item understand the Hardy--Littlewood maximal operator;
    \item state and interpret the maximal theorem;
    \item understand differentiation of integrals;
    \item recognize singular integral operators;
    \item work with the Hilbert transform;
    \item understand Calder\'on--Zygmund kernels and operators;
    \item distinguish strong-type and weak-type estimates;
    \item apply interpolation principles;
    \item understand the Hausdorff--Young inequality;
    \item recognize the roles of Hardy spaces and BMO;
    \item understand Littlewood--Paley frequency decomposition;
    \item interpret harmonic functions through the mean-value property;
    \item work with the Poisson kernel and harmonic extensions;
    \item understand harmonic conjugates and boundary values;
    \item recognize the role of harmonic analysis in PDEs, probability,
          signal processing, and geometric analysis.
\end{itemize}
}

%%%%%%%%%%%%%%%%%%%%%%%%%%%%%%%%%%%%%%%%%%%%%%%%%%%%%%%%%%%%
\subsection{From Fourier Analysis to Harmonic Analysis}
\label{subsec:fourier-to-harmonic}
%%%%%%%%%%%%%%%%%%%%%%%%%%%%%%%%%%%%%%%%%%%%%%%%%%%%%%%%%%%%

Fourier analysis studies the decomposition of functions into frequencies.

Harmonic analysis expands this viewpoint into a broader study of

\[
\boxed{
\text{functions, oscillation, averages, singularities, and operators}.
}
\]

Rather than merely computing Fourier coefficients or transforms,
harmonic analysis asks questions such as

\[
\boxed{
\text{How large can an averaging operator become?}
}
\]

\[
\boxed{
\text{How does an operator behave on }L^p\text{ spaces?}
}
\]

\[
\boxed{
\text{How can frequency information control local behavior?}
}
\]

and

\[
\boxed{
\text{How can local averages recover a function?}
}
\]

The subject is a major bridge between Fourier analysis, functional
analysis, partial differential equations, probability, and geometry.

%%%%%%%%%%%%%%%%%%%%%%%%%%%%%%%%%%%%%%%%%%%%%%%%%%%%%%%%%%%%
\subsection{The Harmonic-Analysis Viewpoint}
\label{subsec:harmonic-analysis-viewpoint}
%%%%%%%%%%%%%%%%%%%%%%%%%%%%%%%%%%%%%%%%%%%%%%%%%%%%%%%%%%%%

A recurring strategy in harmonic analysis is

\[
\boxed{
\text{function}
\longrightarrow
\text{operator}
\longrightarrow
\text{estimate}
\longrightarrow
\text{structural information}.
}
\]

For example, instead of attempting to compute a transformed function
exactly, one may seek an inequality of the form

\[
\boxed{
\|Tf\|_q
\leq
C\|f\|_p.
}
\]

Such estimates reveal whether an operator is stable on a function space.

%%%%%%%%%%%%%%%%%%%%%%%%%%%%%%%%%%%%%%%%%%%%%%%%%%%%%%%%%%%%
\subsection{Translations}
\label{subsec:harmonic-translations}
%%%%%%%%%%%%%%%%%%%%%%%%%%%%%%%%%%%%%%%%%%%%%%%%%%%%%%%%%%%%

For \(y\in\R^n\), define the translation operator

\[
\boxed{
\tau_yf(x)=f(x-y).
}
\]

The lowercase Greek letter \(\tau\) is \emph{tau}.

For \(1\leq p\leq\infty\),

\[
\boxed{
\|\tau_yf\|_p
=
\|f\|_p.
}
\]

Thus translation preserves \(L^p\) norms.

Translation invariance is a central symmetry of Euclidean harmonic
analysis.

%%%%%%%%%%%%%%%%%%%%%%%%%%%%%%%%%%%%%%%%%%%%%%%%%%%%%%%%%%%%
\subsection{Dilations}
\label{subsec:harmonic-dilations}
%%%%%%%%%%%%%%%%%%%%%%%%%%%%%%%%%%%%%%%%%%%%%%%%%%%%%%%%%%%%

For \(t>0\), a basic dilation is

\[
\boxed{
D_tf(x)=f(tx).
}
\]

A change of variables gives

\[
\boxed{
\|D_tf\|_p
=
t^{-n/p}\|f\|_p
}
\]

for \(1\leq p<\infty\).

A normalized \(L^p\)-preserving dilation is therefore

\[
\boxed{
\widetilde{D}_tf(x)
=
t^{n/p}f(tx).
}
\]

Scaling behavior often reveals the correct exponents in analytic
inequalities.

%%%%%%%%%%%%%%%%%%%%%%%%%%%%%%%%%%%%%%%%%%%%%%%%%%%%%%%%%%%%
\subsection{Modulations}
\label{subsec:harmonic-modulations}
%%%%%%%%%%%%%%%%%%%%%%%%%%%%%%%%%%%%%%%%%%%%%%%%%%%%%%%%%%%%

For \(\xi\in\R^n\), define

\[
\boxed{
M_\xi f(x)
=
e^{2\pi i x\cdot\xi}f(x).
}
\]

This operation is called modulation.

Since

\[
|e^{2\pi i x\cdot\xi}|=1,
\]

we have

\[
\boxed{
\|M_\xi f\|_p
=
\|f\|_p.
}
\]

Under Fourier transformation, modulation becomes translation in
frequency space.

%%%%%%%%%%%%%%%%%%%%%%%%%%%%%%%%%%%%%%%%%%%%%%%%%%%%%%%%%%%%
\subsection{Convolution}
\label{subsec:harmonic-convolution}
%%%%%%%%%%%%%%%%%%%%%%%%%%%%%%%%%%%%%%%%%%%%%%%%%%%%%%%%%%%%

For suitable functions \(f,g:\R^n\to\mathbb{C}\), define

\[
\boxed{
(f*g)(x)
=
\int_{\R^n}
f(x-y)g(y)\,dy.
}
\]

Convolution combines translation and averaging.

Important properties include

\[
\boxed{
f*g=g*f,
}
\]

\[
\boxed{
(f*g)*h=f*(g*h),
}
\]

and

\[
\boxed{
f*(g+h)=f*g+f*h.
}
\]

Fourier transformation converts convolution into multiplication:

\[
\boxed{
\widehat{f*g}
=
\widehat{f}\,\widehat{g}.
}
\]

%%%%%%%%%%%%%%%%%%%%%%%%%%%%%%%%%%%%%%%%%%%%%%%%%%%%%%%%%%%%
\subsection{Young's Convolution Inequality}
\label{subsec:young-harmonic}
%%%%%%%%%%%%%%%%%%%%%%%%%%%%%%%%%%%%%%%%%%%%%%%%%%%%%%%%%%%%

\begin{theorem}[Young's Convolution Inequality]
Suppose

\[
1\leq p,q,r\leq\infty
\]

and

\[
\boxed{
1+\frac1r
=
\frac1p+\frac1q.
}
\]

Then, under the standard hypotheses,

\[
\boxed{
\|f*g\|_r
\leq
\|f\|_p\|g\|_q.
}
\]
\end{theorem}

A particularly important case is

\[
\boxed{
\|f*g\|_p
\leq
\|f\|_p\|g\|_1.
}
\]

Thus convolution with an \(L^1\) kernel defines a bounded operator on
\(L^p\).

%%%%%%%%%%%%%%%%%%%%%%%%%%%%%%%%%%%%%%%%%%%%%%%%%%%%%%%%%%%%
\subsection{Approximate Identities}
\label{subsec:harmonic-approximate-identities}
%%%%%%%%%%%%%%%%%%%%%%%%%%%%%%%%%%%%%%%%%%%%%%%%%%%%%%%%%%%%

Let

\[
\phi\in L^1(\R^n)
\]

satisfy

\[
\boxed{
\int_{\R^n}\phi(x)\,dx=1.
}
\]

For \(\varepsilon>0\), define

\[
\boxed{
\phi_\varepsilon(x)
=
\frac{1}{\varepsilon^n}
\phi\left(\frac{x}{\varepsilon}\right).
}
\]

The lowercase Greek letter \(\varepsilon\) is \emph{epsilon}.

As

\[
\varepsilon\to0,
\]

the mass of \(\phi_\varepsilon\) becomes increasingly concentrated near
the origin.

The family

\[
(\phi_\varepsilon)_{\varepsilon>0}
\]

is a basic example of an approximate identity.

%%%%%%%%%%%%%%%%%%%%%%%%%%%%%%%%%%%%%%%%%%%%%%%%%%%%%%%%%%%%
\subsection{Approximation by Convolution}
\label{subsec:approximation-by-convolution}
%%%%%%%%%%%%%%%%%%%%%%%%%%%%%%%%%%%%%%%%%%%%%%%%%%%%%%%%%%%%

Under standard assumptions,

\[
\boxed{
f*\phi_\varepsilon
\to
f
}
\]

as

\[
\varepsilon\to0.
\]

For example, if

\[
1\leq p<\infty
\]

and \(f\in L^p(\R^n)\), then for a standard approximate identity,

\[
\boxed{
\|f*\phi_\varepsilon-f\|_p
\to0.
}
\]

Convolution therefore provides a systematic method of approximating
rough functions by better-behaved functions.

%%%%%%%%%%%%%%%%%%%%%%%%%%%%%%%%%%%%%%%%%%%%%%%%%%%%%%%%%%%%
\subsection{Mollifiers}
\label{subsec:mollifiers}
%%%%%%%%%%%%%%%%%%%%%%%%%%%%%%%%%%%%%%%%%%%%%%%%%%%%%%%%%%%%

A mollifier is typically a smooth, compactly supported approximate
identity.

A standard construction begins with

\[
\phi(x)
=
\begin{cases}
C\exp\left(-\dfrac{1}{1-|x|^2}\right),
& |x|<1,\\[2mm]
0,
& |x|\geq1,
\end{cases}
\]

where \(C\) is chosen so that

\[
\int_{\R^n}\phi(x)\,dx=1.
\]

Then

\[
\boxed{
\phi_\varepsilon(x)
=
\varepsilon^{-n}
\phi(x/\varepsilon).
}
\]

The mollified function is

\[
\boxed{
f_\varepsilon
=
f*\phi_\varepsilon.
}
\]

Under suitable assumptions,

\[
f_\varepsilon\in C^\infty.
\]

%%%%%%%%%%%%%%%%%%%%%%%%%%%%%%%%%%%%%%%%%%%%%%%%%%%%%%%%%%%%
\subsection{Averaging Operators}
\label{subsec:averaging-operators}
%%%%%%%%%%%%%%%%%%%%%%%%%%%%%%%%%%%%%%%%%%%%%%%%%%%%%%%%%%%%

For a ball

\[
B(x,r)\subseteq\R^n,
\]

the average of \(f\) over the ball is

\[
\boxed{
A_rf(x)
=
\frac{1}{|B(x,r)|}
\int_{B(x,r)}f(y)\,dy.
}
\]

Here

\[
|B(x,r)|
\]

denotes Lebesgue measure.

A fundamental question is whether

\[
A_rf(x)
\]

approaches \(f(x)\) as

\[
r\to0.
\]

This leads to the differentiation theory of integrals.

%%%%%%%%%%%%%%%%%%%%%%%%%%%%%%%%%%%%%%%%%%%%%%%%%%%%%%%%%%%%
\subsection{Hardy--Littlewood Maximal Function}
\label{subsec:hardy-littlewood-maximal}
%%%%%%%%%%%%%%%%%%%%%%%%%%%%%%%%%%%%%%%%%%%%%%%%%%%%%%%%%%%%

\begin{definitionbox}{Hardy--Littlewood Maximal Function}
For a locally integrable function \(f\) on \(\R^n\), define

\[
\boxed{
Mf(x)
=
\sup_{r>0}
\frac{1}{|B(x,r)|}
\int_{B(x,r)}
|f(y)|\,dy.
}
\]
\end{definitionbox}

The operator

\[
f\mapsto Mf
\]

is called the Hardy--Littlewood maximal operator.

It measures the largest local average of \(|f|\) around each point.

%%%%%%%%%%%%%%%%%%%%%%%%%%%%%%%%%%%%%%%%%%%%%%%%%%%%%%%%%%%%
\subsection{Why the Maximal Function Matters}
\label{subsec:why-maximal-function}
%%%%%%%%%%%%%%%%%%%%%%%%%%%%%%%%%%%%%%%%%%%%%%%%%%%%%%%%%%%%

The maximal function controls many limiting processes.

For example,

\[
\left|
\frac{1}{|B(x,r)|}
\int_{B(x,r)}
f(y)\,dy
\right|
\leq
Mf(x).
\]

Thus the maximal function acts as a universal pointwise bound for local
averages.

Maximal estimates are used to prove

\[
\boxed{
\text{almost-everywhere convergence}
}
\]

from norm estimates.

%%%%%%%%%%%%%%%%%%%%%%%%%%%%%%%%%%%%%%%%%%%%%%%%%%%%%%%%%%%%
\subsection{Weak-Type Estimates}
\label{subsec:weak-type-estimates}
%%%%%%%%%%%%%%%%%%%%%%%%%%%%%%%%%%%%%%%%%%%%%%%%%%%%%%%%%%%%

\begin{definitionbox}{Weak Type \((p,q)\)}
An operator \(T\) is of weak type \((p,q)\) if there exists \(C>0\)
such that

\[
\boxed{
\mu\left(
\{x:|Tf(x)|>\lambda\}
\right)
\leq
\left(
\frac{C\|f\|_p}{\lambda}
\right)^q
}
\]

for every \(\lambda>0\), whenever the expression is meaningful.
\end{definitionbox}

The lowercase Greek letter \(\lambda\) is \emph{lambda} and represents a
positive threshold.

Weak-type estimates control the size of level sets rather than directly
controlling the \(L^q\) norm.

%%%%%%%%%%%%%%%%%%%%%%%%%%%%%%%%%%%%%%%%%%%%%%%%%%%%%%%%%%%%
\subsection{Strong-Type Estimates}
\label{subsec:strong-type-estimates}
%%%%%%%%%%%%%%%%%%%%%%%%%%%%%%%%%%%%%%%%%%%%%%%%%%%%%%%%%%%%

\begin{definitionbox}{Strong Type \((p,q)\)}
An operator \(T\) is of strong type \((p,q)\) if

\[
\boxed{
\|Tf\|_q
\leq
C\|f\|_p
}
\]

for some constant \(C\) independent of \(f\).
\end{definitionbox}

Strong-type estimates are ordinary boundedness statements between
\(L^p\) spaces.

%%%%%%%%%%%%%%%%%%%%%%%%%%%%%%%%%%%%%%%%%%%%%%%%%%%%%%%%%%%%
\subsection{Hardy--Littlewood Maximal Theorem}
\label{subsec:hardy-littlewood-maximal-theorem}
%%%%%%%%%%%%%%%%%%%%%%%%%%%%%%%%%%%%%%%%%%%%%%%%%%%%%%%%%%%%

\begin{theorem}[Hardy--Littlewood Maximal Theorem]
The Hardy--Littlewood maximal operator is of weak type \((1,1)\):

\[
\boxed{
\left|
\{x\in\R^n:Mf(x)>\lambda\}
\right|
\leq
\frac{C_n}{\lambda}
\|f\|_1.
}
\]

Moreover, for

\[
1<p\leq\infty,
\]

the maximal operator is bounded on \(L^p\):

\[
\boxed{
\|Mf\|_p
\leq
C_{n,p}\|f\|_p.
}
\]
\end{theorem}

The constants depend on the dimension and, in the strong estimate, on
\(p\).

%%%%%%%%%%%%%%%%%%%%%%%%%%%%%%%%%%%%%%%%%%%%%%%%%%%%%%%%%%%%
\subsection{Covering Lemmas}
\label{subsec:covering-lemmas}
%%%%%%%%%%%%%%%%%%%%%%%%%%%%%%%%%%%%%%%%%%%%%%%%%%%%%%%%%%%%

Maximal-function estimates rely on geometric covering arguments.

A typical problem begins with a family of balls

\[
\{B_\alpha\}
\]

covering a set.

A covering lemma extracts a controlled subcollection of disjoint or
nearly disjoint balls.

The Vitali covering principle and related covering lemmas are fundamental
tools in real-variable harmonic analysis.

%%%%%%%%%%%%%%%%%%%%%%%%%%%%%%%%%%%%%%%%%%%%%%%%%%%%%%%%%%%%
\subsection{Lebesgue Differentiation Theorem}
\label{subsec:lebesgue-differentiation}
%%%%%%%%%%%%%%%%%%%%%%%%%%%%%%%%%%%%%%%%%%%%%%%%%%%%%%%%%%%%

\begin{theorem}[Lebesgue Differentiation Theorem]
If

\[
f\in L^1_{\mathrm{loc}}(\R^n),
\]

then for almost every \(x\in\R^n\),

\[
\boxed{
\lim_{r\to0}
\frac{1}{|B(x,r)|}
\int_{B(x,r)}
f(y)\,dy
=
f(x).
}
\]
\end{theorem}

Equivalently, for almost every \(x\),

\[
\boxed{
\lim_{r\to0}
\frac{1}{|B(x,r)|}
\int_{B(x,r)}
|f(y)-f(x)|\,dy
=
0.
}
\]

Thus an integrable function can be recovered almost everywhere from its
local averages.

%%%%%%%%%%%%%%%%%%%%%%%%%%%%%%%%%%%%%%%%%%%%%%%%%%%%%%%%%%%%
\subsection{Lebesgue Points}
\label{subsec:lebesgue-points}
%%%%%%%%%%%%%%%%%%%%%%%%%%%%%%%%%%%%%%%%%%%%%%%%%%%%%%%%%%%%

\begin{definitionbox}{Lebesgue Point}
A point \(x\) is a Lebesgue point of \(f\) if

\[
\boxed{
\lim_{r\to0}
\frac{1}{|B(x,r)|}
\int_{B(x,r)}
|f(y)-f(x)|\,dy
=
0.
}
\]
\end{definitionbox}

The Lebesgue Differentiation Theorem states that almost every point is a
Lebesgue point of a locally integrable function.

%%%%%%%%%%%%%%%%%%%%%%%%%%%%%%%%%%%%%%%%%%%%%%%%%%%%%%%%%%%%
\subsection{Interpolation}
\label{subsec:harmonic-interpolation}
%%%%%%%%%%%%%%%%%%%%%%%%%%%%%%%%%%%%%%%%%%%%%%%%%%%%%%%%%%%%

Suppose an operator is bounded between certain endpoint spaces.

Interpolation asks whether boundedness also holds for intermediate
exponents.

Schematically,

\[
\boxed{
L^{p_0}\to L^{q_0}
}
\]

and

\[
\boxed{
L^{p_1}\to L^{q_1}
}
\]

may imply

\[
\boxed{
L^p\to L^q
}
\]

for intermediate \(p,q\).

This principle is extremely powerful because endpoint estimates are often
easier to establish.

%%%%%%%%%%%%%%%%%%%%%%%%%%%%%%%%%%%%%%%%%%%%%%%%%%%%%%%%%%%%
\subsection{Riesz--Thorin Interpolation Theorem}
\label{subsec:riesz-thorin}
%%%%%%%%%%%%%%%%%%%%%%%%%%%%%%%%%%%%%%%%%%%%%%%%%%%%%%%%%%%%

\begin{theorem}[Riesz--Thorin Interpolation Theorem]
Suppose \(T\) is linear and bounded

\[
T:L^{p_0}\to L^{q_0}
\]

and

\[
T:L^{p_1}\to L^{q_1}.
\]

For \(0<\theta<1\), define

\[
\boxed{
\frac1{p_\theta}
=
\frac{1-\theta}{p_0}
+
\frac{\theta}{p_1}
}
\]

and

\[
\boxed{
\frac1{q_\theta}
=
\frac{1-\theta}{q_0}
+
\frac{\theta}{q_1}.
}
\]

Then

\[
\boxed{
T:L^{p_\theta}\to L^{q_\theta}
}
\]

is bounded, with the corresponding interpolated norm estimate.
\end{theorem}

The lowercase Greek letter \(\theta\) is \emph{theta} and represents the
interpolation parameter.

%%%%%%%%%%%%%%%%%%%%%%%%%%%%%%%%%%%%%%%%%%%%%%%%%%%%%%%%%%%%
\subsection{Marcinkiewicz Interpolation}
\label{subsec:marcinkiewicz-interpolation}
%%%%%%%%%%%%%%%%%%%%%%%%%%%%%%%%%%%%%%%%%%%%%%%%%%%%%%%%%%%%

The Marcinkiewicz Interpolation Theorem is designed for sublinear
operators and weak-type endpoint estimates.

In simplified form, if a sublinear operator has suitable weak-type
bounds at two endpoints, then it has strong-type bounds for intermediate
exponents.

This theorem is one route from the weak \((1,1)\) maximal estimate to

\[
\boxed{
\|Mf\|_p
\leq
C_p\|f\|_p,
\qquad
1<p<\infty.
}
\]

%%%%%%%%%%%%%%%%%%%%%%%%%%%%%%%%%%%%%%%%%%%%%%%%%%%%%%%%%%%%
\subsection{Hausdorff--Young Inequality}
\label{subsec:hausdorff-young}
%%%%%%%%%%%%%%%%%%%%%%%%%%%%%%%%%%%%%%%%%%%%%%%%%%%%%%%%%%%%

Plancherel gives

\[
L^2\to L^2
\]

boundedness of the Fourier transform.

The elementary estimate

\[
\boxed{
\|\widehat{f}\|_\infty
\leq
\|f\|_1
}
\]

gives

\[
L^1\to L^\infty.
\]

Interpolation between these endpoints yields the Hausdorff--Young
inequality.

\begin{theorem}[Hausdorff--Young Inequality]
Let

\[
1\leq p\leq2
\]

and let \(p'\) satisfy

\[
\boxed{
\frac1p+\frac1{p'}=1.
}
\]

Then

\[
\boxed{
\|\widehat{f}\|_{p'}
\leq
C_p\|f\|_p.
}
\]
\end{theorem}

Under common Fourier normalizations, the non-sharp interpolation argument
may be expressed with \(C_p\leq1\).

%%%%%%%%%%%%%%%%%%%%%%%%%%%%%%%%%%%%%%%%%%%%%%%%%%%%%%%%%%%%
\subsection{Singular Integrals}
\label{subsec:singular-integrals}
%%%%%%%%%%%%%%%%%%%%%%%%%%%%%%%%%%%%%%%%%%%%%%%%%%%%%%%%%%%%

Convolution with an integrable kernel is relatively well behaved.

Harmonic analysis also studies kernels that are singular at the origin.

Formally,

\[
\boxed{
Tf(x)
=
\int_{\R^n}
K(x-y)f(y)\,dy,
}
\]

where \(K\) may fail to be integrable near

\[
x=y.
\]

Such integrals require cancellation or principal-value interpretation.

%%%%%%%%%%%%%%%%%%%%%%%%%%%%%%%%%%%%%%%%%%%%%%%%%%%%%%%%%%%%
\subsection{Principal Value Integrals}
\label{subsec:principal-value-integrals}
%%%%%%%%%%%%%%%%%%%%%%%%%%%%%%%%%%%%%%%%%%%%%%%%%%%%%%%%%%%%

A singular integral may be interpreted through a Cauchy principal value.

For example,

\[
\boxed{
\operatorname{p.v.}
\int_{\R}
\frac{f(y)}{x-y}\,dy
=
\lim_{\varepsilon\to0^+}
\int_{|x-y|>\varepsilon}
\frac{f(y)}{x-y}\,dy,
}
\]

provided the limit exists.

The notation

\[
\operatorname{p.v.}
\]

means \emph{principal value}.

%%%%%%%%%%%%%%%%%%%%%%%%%%%%%%%%%%%%%%%%%%%%%%%%%%%%%%%%%%%%
\subsection{Hilbert Transform}
\label{subsec:hilbert-transform}
%%%%%%%%%%%%%%%%%%%%%%%%%%%%%%%%%%%%%%%%%%%%%%%%%%%%%%%%%%%%

\begin{definitionbox}{Hilbert Transform}
For a sufficiently regular function \(f\), the Hilbert transform is

\[
\boxed{
Hf(x)
=
\frac1\pi
\operatorname{p.v.}
\int_{\R}
\frac{f(y)}{x-y}\,dy.
}
\]
\end{definitionbox}

The Hilbert transform is one of the fundamental singular integral
operators.

It may be viewed formally as convolution with

\[
\boxed{
\frac{1}{\pi x},
}
\]

interpreted in the principal-value sense.

%%%%%%%%%%%%%%%%%%%%%%%%%%%%%%%%%%%%%%%%%%%%%%%%%%%%%%%%%%%%
\subsection{Hilbert Transform in Frequency Space}
\label{subsec:hilbert-transform-frequency}
%%%%%%%%%%%%%%%%%%%%%%%%%%%%%%%%%%%%%%%%%%%%%%%%%%%%%%%%%%%%

Under the Fourier convention

\[
\widehat{f}(\xi)
=
\int_{\R}
f(x)e^{-2\pi ix\xi}\,dx,
\]

the Hilbert transform satisfies

\[
\boxed{
\widehat{Hf}(\xi)
=
-i\,\operatorname{sgn}(\xi)\widehat{f}(\xi).
}
\]

Here

\[
\operatorname{sgn}(\xi)
=
\begin{cases}
1, & \xi>0,\\
0, & \xi=0,\\
-1, & \xi<0.
\end{cases}
\]

Thus the Hilbert transform is a Fourier multiplier operator.

%%%%%%%%%%%%%%%%%%%%%%%%%%%%%%%%%%%%%%%%%%%%%%%%%%%%%%%%%%%%
\subsection{\(L^p\) Boundedness of the Hilbert Transform}
\label{subsec:hilbert-lp-boundedness}
%%%%%%%%%%%%%%%%%%%%%%%%%%%%%%%%%%%%%%%%%%%%%%%%%%%%%%%%%%%%

A fundamental theorem states that

\[
\boxed{
\|Hf\|_p
\leq
C_p\|f\|_p
}
\]

for

\[
\boxed{
1<p<\infty.
}
\]

At \(p=1\), strong \(L^1\) boundedness fails, but a weak-type estimate is
available.

This endpoint behavior is typical of singular integral operators.

%%%%%%%%%%%%%%%%%%%%%%%%%%%%%%%%%%%%%%%%%%%%%%%%%%%%%%%%%%%%
\subsection{Calder\'on--Zygmund Kernels}
\label{subsec:calderon-zygmund-kernels}
%%%%%%%%%%%%%%%%%%%%%%%%%%%%%%%%%%%%%%%%%%%%%%%%%%%%%%%%%%%%

A Calder\'on--Zygmund kernel is a singular kernel satisfying appropriate
size, smoothness, and cancellation conditions.

A typical size estimate is

\[
\boxed{
|K(x)|
\leq
\frac{C}{|x|^n},
\qquad
x\neq0.
}
\]

A typical smoothness estimate has the form

\[
\boxed{
|K(x-h)-K(x)|
\leq
C\frac{|h|^\alpha}{|x|^{n+\alpha}}
}
\]

when

\[
|h|
\leq
\frac{|x|}{2},
\]

for some

\[
0<\alpha\leq1.
\]

The lowercase Greek letter \(\alpha\) is \emph{alpha} and here measures
kernel regularity.

%%%%%%%%%%%%%%%%%%%%%%%%%%%%%%%%%%%%%%%%%%%%%%%%%%%%%%%%%%%%
\subsection{Calder\'on--Zygmund Operators}
\label{subsec:calderon-zygmund-operators}
%%%%%%%%%%%%%%%%%%%%%%%%%%%%%%%%%%%%%%%%%%%%%%%%%%%%%%%%%%%%

Very roughly, a Calder\'on--Zygmund operator is a singular integral
operator associated with a Calder\'on--Zygmund kernel and satisfying an
appropriate \(L^2\) boundedness condition.

A central result of the theory is that such operators extend boundedly to

\[
\boxed{
L^p(\R^n),
\qquad
1<p<\infty.
}
\]

They also satisfy important endpoint estimates.

The Hilbert transform is the basic one-dimensional example.

%%%%%%%%%%%%%%%%%%%%%%%%%%%%%%%%%%%%%%%%%%%%%%%%%%%%%%%%%%%%
\subsection{Riesz Transforms}
\label{subsec:riesz-transforms}
%%%%%%%%%%%%%%%%%%%%%%%%%%%%%%%%%%%%%%%%%%%%%%%%%%%%%%%%%%%%

In \(\R^n\), the Riesz transforms generalize the Hilbert transform.

Their Fourier multipliers are

\[
\boxed{
\widehat{R_jf}(\xi)
=
-i
\frac{\xi_j}{|\xi|}
\widehat{f}(\xi),
\qquad
j=1,\ldots,n.
}
\]

These operators are bounded on

\[
\boxed{
L^p(\R^n),
\qquad
1<p<\infty.
}
\]

Riesz transforms are important in PDEs and potential theory.

%%%%%%%%%%%%%%%%%%%%%%%%%%%%%%%%%%%%%%%%%%%%%%%%%%%%%%%%%%%%
\subsection{Fourier Multipliers}
\label{subsec:fourier-multipliers}
%%%%%%%%%%%%%%%%%%%%%%%%%%%%%%%%%%%%%%%%%%%%%%%%%%%%%%%%%%%%

Suppose \(m:\R^n\to\mathbb{C}\) is a function.

A Fourier multiplier operator is formally defined by

\[
\boxed{
\widehat{Tf}(\xi)
=
m(\xi)\widehat{f}(\xi).
}
\]

The function \(m\) is called the multiplier or symbol.

By Plancherel's Theorem, if

\[
m\in L^\infty,
\]

then

\[
\boxed{
\|Tf\|_2
\leq
\|m\|_\infty\|f\|_2.
}
\]

The difficult question is often whether \(T\) is bounded on \(L^p\) for
\(p\neq2\).

%%%%%%%%%%%%%%%%%%%%%%%%%%%%%%%%%%%%%%%%%%%%%%%%%%%%%%%%%%%%
\subsection{Multiplier Theorems}
\label{subsec:multiplier-theorems}
%%%%%%%%%%%%%%%%%%%%%%%%%%%%%%%%%%%%%%%%%%%%%%%%%%%%%%%%%%%%

Multiplier theorems give conditions on \(m\) ensuring that

\[
\boxed{
T_m:L^p\to L^p
}
\]

is bounded.

Typically, smoothness and decay of derivatives of \(m\) compensate for
the lack of Hilbert-space orthogonality when

\[
p\neq2.
\]

The Mikhlin multiplier theorem is a major example and is developed more
fully in advanced harmonic analysis and PDE theory.

%%%%%%%%%%%%%%%%%%%%%%%%%%%%%%%%%%%%%%%%%%%%%%%%%%%%%%%%%%%%
\subsection{Frequency Localization}
\label{subsec:frequency-localization}
%%%%%%%%%%%%%%%%%%%%%%%%%%%%%%%%%%%%%%%%%%%%%%%%%%%%%%%%%%%%

Instead of studying all frequencies simultaneously, harmonic analysis
often decomposes a function into frequency bands.

One may write schematically

\[
\boxed{
f
=
\sum_{j\in\Z}
\Delta_jf,
}
\]

where \(\Delta_jf\) contains frequencies roughly satisfying

\[
\boxed{
|\xi|\approx2^j.
}
\]

This is called a dyadic frequency decomposition.

The word \emph{dyadic} refers to powers of \(2\).

%%%%%%%%%%%%%%%%%%%%%%%%%%%%%%%%%%%%%%%%%%%%%%%%%%%%%%%%%%%%
\subsection{Littlewood--Paley Decomposition}
\label{subsec:littlewood-paley}
%%%%%%%%%%%%%%%%%%%%%%%%%%%%%%%%%%%%%%%%%%%%%%%%%%%%%%%%%%%%

Choose smooth functions in frequency space that partition the nonzero
frequencies into overlapping dyadic annuli.

The corresponding operators

\[
\Delta_j
\]

localize \(f\) to frequencies of size approximately \(2^j\).

Schematically,

\[
\boxed{
\widehat{\Delta_jf}(\xi)
=
\psi(2^{-j}\xi)\widehat{f}(\xi),
}
\]

where \(\psi\) is supported away from the origin.

Then

\[
\boxed{
f
\sim
\sum_j\Delta_jf.
}
\]

%%%%%%%%%%%%%%%%%%%%%%%%%%%%%%%%%%%%%%%%%%%%%%%%%%%%%%%%%%%%
\subsection{Littlewood--Paley Square Function}
\label{subsec:littlewood-paley-square}
%%%%%%%%%%%%%%%%%%%%%%%%%%%%%%%%%%%%%%%%%%%%%%%%%%%%%%%%%%%%

The Littlewood--Paley square function is schematically

\[
\boxed{
S(f)(x)
=
\left(
\sum_{j\in\Z}
|\Delta_jf(x)|^2
\right)^{1/2}.
}
\]

A central principle is that, for

\[
1<p<\infty,
\]

the size of \(f\) in \(L^p\) is comparable to the size of its square
function:

\[
\boxed{
\|S(f)\|_p
\asymp
\|f\|_p.
}
\]

The notation

\[
A\asymp B
\]

means that each quantity is bounded by a constant multiple of the other.

%%%%%%%%%%%%%%%%%%%%%%%%%%%%%%%%%%%%%%%%%%%%%%%%%%%%%%%%%%%%
\subsection{Why Dyadic Decomposition Works}
\label{subsec:why-dyadic}
%%%%%%%%%%%%%%%%%%%%%%%%%%%%%%%%%%%%%%%%%%%%%%%%%%%%%%%%%%%%

Frequency localization allows complicated operators to be studied one
scale at a time.

Instead of analyzing

\[
Tf
\]

globally, one analyzes

\[
T(\Delta_jf)
\]

for each frequency scale.

This separates phenomena according to

\[
\boxed{
\text{low frequency},
\qquad
\text{intermediate frequency},
\qquad
\text{high frequency}.
}
\]

The method is central to modern PDE theory.

%%%%%%%%%%%%%%%%%%%%%%%%%%%%%%%%%%%%%%%%%%%%%%%%%%%%%%%%%%%%
\subsection{Hardy Spaces}
\label{subsec:hardy-spaces}
%%%%%%%%%%%%%%%%%%%%%%%%%%%%%%%%%%%%%%%%%%%%%%%%%%%%%%%%%%%%

Classical \(L^p\) spaces become more difficult at or below the endpoint

\[
p=1.
\]

Hardy spaces provide a replacement with better harmonic-analytic
properties.

The real Hardy space

\[
\boxed{
H^1(\R^n)
}
\]

can be characterized in several equivalent ways, including through

\[
\boxed{
\text{maximal functions},
\quad
\text{atomic decompositions},
\quad
\text{square functions}.
}
\]

Hardy spaces allow singular integral theory to extend beyond ordinary
\(L^1\) estimates.

%%%%%%%%%%%%%%%%%%%%%%%%%%%%%%%%%%%%%%%%%%%%%%%%%%%%%%%%%%%%
\subsection{Atomic Decomposition}
\label{subsec:hardy-atomic}
%%%%%%%%%%%%%%%%%%%%%%%%%%%%%%%%%%%%%%%%%%%%%%%%%%%%%%%%%%%%

An \(H^1\) atom is, roughly, a localized function \(a\) satisfying

\[
\boxed{
\operatorname{supp}(a)\subseteq B,
}
\]

\[
\boxed{
\|a\|_\infty
\leq
\frac{1}{|B|},
}
\]

and the cancellation condition

\[
\boxed{
\int a(x)\,dx=0.
}
\]

Functions in \(H^1\) can be decomposed schematically as

\[
\boxed{
f
=
\sum_j\lambda_ja_j
}
\]

with

\[
\boxed{
\sum_j|\lambda_j|<\infty.
}
\]

This reduces many estimates to estimates on simple localized atoms.

%%%%%%%%%%%%%%%%%%%%%%%%%%%%%%%%%%%%%%%%%%%%%%%%%%%%%%%%%%%%
\subsection{BMO}
\label{subsec:bmo}
%%%%%%%%%%%%%%%%%%%%%%%%%%%%%%%%%%%%%%%%%%%%%%%%%%%%%%%%%%%%

\begin{definitionbox}{Bounded Mean Oscillation}
A locally integrable function \(f\) belongs to

\[
\operatorname{BMO}(\R^n)
\]

if

\[
\boxed{
\|f\|_{\operatorname{BMO}}
=
\sup_B
\frac{1}{|B|}
\int_B
|f(x)-f_B|\,dx
<\infty,
}
\]

where

\[
\boxed{
f_B
=
\frac{1}{|B|}
\int_Bf(x)\,dx.
}
\]
\end{definitionbox}

BMO stands for \emph{bounded mean oscillation}.

The supremum is taken over all balls \(B\subseteq\R^n\).

%%%%%%%%%%%%%%%%%%%%%%%%%%%%%%%%%%%%%%%%%%%%%%%%%%%%%%%%%%%%
\subsection{BMO Versus \(L^\infty\)}
\label{subsec:bmo-versus-linfty}
%%%%%%%%%%%%%%%%%%%%%%%%%%%%%%%%%%%%%%%%%%%%%%%%%%%%%%%%%%%%

Every essentially bounded function belongs to BMO:

\[
\boxed{
L^\infty
\subseteq
\operatorname{BMO}.
}
\]

However, BMO contains unbounded functions.

Thus

\[
\boxed{
L^\infty
\subsetneq
\operatorname{BMO}.
}
\]

BMO measures oscillation rather than absolute size.

%%%%%%%%%%%%%%%%%%%%%%%%%%%%%%%%%%%%%%%%%%%%%%%%%%%%%%%%%%%%
\subsection{John--Nirenberg Phenomenon}
\label{subsec:john-nirenberg}
%%%%%%%%%%%%%%%%%%%%%%%%%%%%%%%%%%%%%%%%%%%%%%%%%%%%%%%%%%%%

The John--Nirenberg inequality states that BMO functions have
exponentially decaying oscillation tails.

Very roughly,

\[
\boxed{
\frac{
|\{x\in B:|f(x)-f_B|>\lambda\}|
}{
|B|
}
\leq
C
e^{-c\lambda/\|f\|_{\operatorname{BMO}}}.
}
\]

Thus bounded mean oscillation is much stronger than the definition alone
might suggest.

%%%%%%%%%%%%%%%%%%%%%%%%%%%%%%%%%%%%%%%%%%%%%%%%%%%%%%%%%%%%
\subsection{Hardy--BMO Duality}
\label{subsec:hardy-bmo-duality}
%%%%%%%%%%%%%%%%%%%%%%%%%%%%%%%%%%%%%%%%%%%%%%%%%%%%%%%%%%%%

A fundamental result is

\[
\boxed{
(H^1)^*
\cong
\operatorname{BMO},
}
\]

with appropriate identification modulo constants on the BMO side.

This relationship is analogous to familiar \(L^p\) duality, but it occurs
at the endpoint where ordinary \(L^1\)-\(L^\infty\) theory is often too
rigid for singular integral problems.

%%%%%%%%%%%%%%%%%%%%%%%%%%%%%%%%%%%%%%%%%%%%%%%%%%%%%%%%%%%%
\subsection{Harmonic Functions}
\label{subsec:harmonic-functions}
%%%%%%%%%%%%%%%%%%%%%%%%%%%%%%%%%%%%%%%%%%%%%%%%%%%%%%%%%%%%

Harmonic analysis takes its name partly from the study of harmonic
functions.

\begin{definitionbox}{Harmonic Function}
A twice continuously differentiable function

\[
u:\Omega\to\R
\]

is harmonic if

\[
\boxed{
\Delta u=0
}
\]

throughout the open set \(\Omega\subseteq\R^n\).
\end{definitionbox}

The Laplacian is

\[
\boxed{
\Delta
=
\frac{\partial^2}{\partial x_1^2}
+\cdots+
\frac{\partial^2}{\partial x_n^2}.
}
\]

%%%%%%%%%%%%%%%%%%%%%%%%%%%%%%%%%%%%%%%%%%%%%%%%%%%%%%%%%%%%
\subsection{Mean-Value Property}
\label{subsec:harmonic-mean-value}
%%%%%%%%%%%%%%%%%%%%%%%%%%%%%%%%%%%%%%%%%%%%%%%%%%%%%%%%%%%%

A remarkable property of harmonic functions is that their value at the
center of a ball equals their average over the ball.

If

\[
\overline{B(x,r)}
\subseteq\Omega,
\]

then

\[
\boxed{
u(x)
=
\frac{1}{|B(x,r)|}
\int_{B(x,r)}u(y)\,dy.
}
\]

There is also a corresponding average over the boundary sphere.

Thus harmonic functions are perfectly balanced with respect to local
averaging.

%%%%%%%%%%%%%%%%%%%%%%%%%%%%%%%%%%%%%%%%%%%%%%%%%%%%%%%%%%%%
\subsection{Maximum Principle}
\label{subsec:harmonic-maximum-principle}
%%%%%%%%%%%%%%%%%%%%%%%%%%%%%%%%%%%%%%%%%%%%%%%%%%%%%%%%%%%%

\begin{theorem}[Maximum Principle]
Let \(\Omega\) be a bounded connected domain and let

\[
u\in C^2(\Omega)\cap C(\overline{\Omega})
\]

be harmonic.

Then

\[
\boxed{
\max_{\overline{\Omega}}u
=
\max_{\partial\Omega}u
}
\]

and

\[
\boxed{
\min_{\overline{\Omega}}u
=
\min_{\partial\Omega}u.
}
\]
\end{theorem}

Unless \(u\) is constant, interior points cannot contain a strict global
maximum or minimum.

%%%%%%%%%%%%%%%%%%%%%%%%%%%%%%%%%%%%%%%%%%%%%%%%%%%%%%%%%%%%
\subsection{Poisson Kernel for the Unit Disk}
\label{subsec:poisson-kernel-disk}
%%%%%%%%%%%%%%%%%%%%%%%%%%%%%%%%%%%%%%%%%%%%%%%%%%%%%%%%%%%%

For the unit disk, the Poisson kernel is

\[
\boxed{
P_r(\theta)
=
\frac{1-r^2}
{
1-2r\cos\theta+r^2
},
\qquad
0\leq r<1.
}
\]

If boundary data is given by \(f\), its harmonic extension is

\[
\boxed{
u(re^{i\theta})
=
\frac{1}{2\pi}
\int_{-\pi}^{\pi}
P_r(\theta-t)f(t)\,dt.
}
\]

Thus harmonic extension is a convolution-type averaging process.

%%%%%%%%%%%%%%%%%%%%%%%%%%%%%%%%%%%%%%%%%%%%%%%%%%%%%%%%%%%%
\subsection{Poisson Kernel as an Approximate Identity}
\label{subsec:poisson-approximate-identity}
%%%%%%%%%%%%%%%%%%%%%%%%%%%%%%%%%%%%%%%%%%%%%%%%%%%%%%%%%%%%

As

\[
r\to1^-,
\]

the Poisson kernel becomes increasingly concentrated near

\[
\theta=0.
\]

Also,

\[
\boxed{
\frac{1}{2\pi}
\int_{-\pi}^{\pi}
P_r(\theta)\,d\theta
=
1.
}
\]

Thus the family

\[
P_r
\]

acts as an approximate identity on the circle.

Under suitable assumptions,

\[
\boxed{
P_r*f
\to
f
}
\]

as \(r\to1^-\).

%%%%%%%%%%%%%%%%%%%%%%%%%%%%%%%%%%%%%%%%%%%%%%%%%%%%%%%%%%%%
\subsection{Harmonic Extension of Fourier Series}
\label{subsec:harmonic-extension-fourier}
%%%%%%%%%%%%%%%%%%%%%%%%%%%%%%%%%%%%%%%%%%%%%%%%%%%%%%%%%%%%

Suppose

\[
f(\theta)
\sim
\sum_{n\in\Z}
c_ne^{in\theta}.
\]

Its Poisson extension is

\[
\boxed{
u(r,\theta)
=
\sum_{n\in\Z}
r^{|n|}
c_ne^{in\theta}.
}
\]

The factor

\[
r^{|n|}
\]

damps high frequencies.

For each

\[
r<1,
\]

the resulting function is smoother than the original boundary data.

%%%%%%%%%%%%%%%%%%%%%%%%%%%%%%%%%%%%%%%%%%%%%%%%%%%%%%%%%%%%
\subsection{Harmonic Conjugates}
\label{subsec:harmonic-conjugates}
%%%%%%%%%%%%%%%%%%%%%%%%%%%%%%%%%%%%%%%%%%%%%%%%%%%%%%%%%%%%

In two dimensions, suppose

\[
F(z)=u(x,y)+iv(x,y)
\]

is holomorphic.

Then \(u\) and \(v\) satisfy the Cauchy--Riemann equations

\[
\boxed{
u_x=v_y,
\qquad
u_y=-v_x.
}
\]

Both \(u\) and \(v\) are harmonic.

The function \(v\) is called a harmonic conjugate of \(u\).

Boundary values of harmonic conjugates are closely connected to the
Hilbert transform.

%%%%%%%%%%%%%%%%%%%%%%%%%%%%%%%%%%%%%%%%%%%%%%%%%%%%%%%%%%%%
\subsection{Subharmonic Functions}
\label{subsec:subharmonic-functions}
%%%%%%%%%%%%%%%%%%%%%%%%%%%%%%%%%%%%%%%%%%%%%%%%%%%%%%%%%%%%

A sufficiently smooth function \(u\) is subharmonic when

\[
\boxed{
\Delta u\geq0.
}
\]

A subharmonic function satisfies a mean-value inequality rather than a
mean-value equality.

Roughly,

\[
\boxed{
u(x)
\leq
\text{average of }u\text{ around }x.
}
\]

Subharmonic functions are important in complex analysis, potential
theory, and PDEs.

%%%%%%%%%%%%%%%%%%%%%%%%%%%%%%%%%%%%%%%%%%%%%%%%%%%%%%%%%%%%
\subsection{Potential-Theoretic Viewpoint}
\label{subsec:harmonic-potential-viewpoint}
%%%%%%%%%%%%%%%%%%%%%%%%%%%%%%%%%%%%%%%%%%%%%%%%%%%%%%%%%%%%

The Laplace equation

\[
\Delta u=0
\]

models equilibrium phenomena.

The Poisson equation

\[
\boxed{
-\Delta u=f
}
\]

models a field generated by a source \(f\).

Harmonic analysis studies these equations using

\[
\boxed{
\text{Fourier transforms},
\quad
\text{singular integrals},
\quad
\text{maximal functions},
\quad
\text{function-space estimates}.
}
\]

This provides a direct bridge to potential theory and PDEs.

%%%%%%%%%%%%%%%%%%%%%%%%%%%%%%%%%%%%%%%%%%%%%%%%%%%%%%%%%%%%
\subsection{Fractional Integration}
\label{subsec:fractional-integration}
%%%%%%%%%%%%%%%%%%%%%%%%%%%%%%%%%%%%%%%%%%%%%%%%%%%%%%%%%%%%

Harmonic analysis also studies integral operators whose kernels have
power-law singularities.

A model operator is the Riesz potential

\[
\boxed{
I_\alpha f(x)
=
C_{n,\alpha}
\int_{\R^n}
\frac{f(y)}
{|x-y|^{n-\alpha}}
\,dy,
}
\]

where

\[
0<\alpha<n.
\]

These operators behave like negative fractional powers of the Laplacian.

%%%%%%%%%%%%%%%%%%%%%%%%%%%%%%%%%%%%%%%%%%%%%%%%%%%%%%%%%%%%
\subsection{Hardy--Littlewood--Sobolev Principle}
\label{subsec:hls-principle}
%%%%%%%%%%%%%%%%%%%%%%%%%%%%%%%%%%%%%%%%%%%%%%%%%%%%%%%%%%%%

Fractional integration improves integrability.

Under suitable exponent conditions,

\[
\boxed{
I_\alpha:L^p(\R^n)\to L^q(\R^n),
}
\]

where

\[
\boxed{
\frac1q
=
\frac1p-\frac{\alpha}{n}.
}
\]

This relationship is a central example of how scaling predicts the
correct exponents in an analytic inequality.

%%%%%%%%%%%%%%%%%%%%%%%%%%%%%%%%%%%%%%%%%%%%%%%%%%%%%%%%%%%%
\subsection{Scaling Analysis}
\label{subsec:harmonic-scaling-analysis}
%%%%%%%%%%%%%%%%%%%%%%%%%%%%%%%%%%%%%%%%%%%%%%%%%%%%%%%%%%%%

Suppose an inequality is proposed:

\[
\boxed{
\|Tf\|_q
\leq
C\|f\|_p.
}
\]

A useful first test is to replace \(f\) by a rescaled function

\[
f_\lambda(x)=f(\lambda x).
\]

If the two sides scale differently, the inequality cannot hold uniformly
for all \(\lambda>0\).

Scaling therefore provides a fast way to identify necessary relationships
between exponents.

%%%%%%%%%%%%%%%%%%%%%%%%%%%%%%%%%%%%%%%%%%%%%%%%%%%%%%%%%%%%
\subsection{Oscillation and Cancellation}
\label{subsec:oscillation-cancellation}
%%%%%%%%%%%%%%%%%%%%%%%%%%%%%%%%%%%%%%%%%%%%%%%%%%%%%%%%%%%%

Absolute values often destroy useful information.

For example,

\[
\int
e^{i\lambda\phi(x)}
a(x)\,dx
\]

may be small even when

\[
\int|a(x)|\,dx
\]

is large.

The reduction occurs because positive and negative complex phases cancel.

Harmonic analysis therefore studies not only magnitude but also

\[
\boxed{
\text{oscillation and cancellation}.
}
\]

These ideas lead to oscillatory integral estimates and stationary-phase
methods.

%%%%%%%%%%%%%%%%%%%%%%%%%%%%%%%%%%%%%%%%%%%%%%%%%%%%%%%%%%%%
\subsection{Oscillatory Integrals}
\label{subsec:oscillatory-integrals}
%%%%%%%%%%%%%%%%%%%%%%%%%%%%%%%%%%%%%%%%%%%%%%%%%%%%%%%%%%%%

A typical oscillatory integral has the form

\[
\boxed{
I(\lambda)
=
\int_{\R^n}
e^{i\lambda\phi(x)}
a(x)\,dx.
}
\]

Here

\[
\lambda
\]

is a large parameter,

\[
\phi
\]

is the phase function, and

\[
a
\]

is the amplitude.

One seeks decay estimates such as

\[
\boxed{
|I(\lambda)|
\leq
C\lambda^{-\beta}.
}
\]

The geometry of the phase determines the decay.

%%%%%%%%%%%%%%%%%%%%%%%%%%%%%%%%%%%%%%%%%%%%%%%%%%%%%%%%%%%%
\subsection{Stationary Phase Principle}
\label{subsec:stationary-phase-preview}
%%%%%%%%%%%%%%%%%%%%%%%%%%%%%%%%%%%%%%%%%%%%%%%%%%%%%%%%%%%%

Rapid oscillation generally produces cancellation except near points
where the phase changes slowly.

These are stationary points satisfying

\[
\boxed{
\nabla\phi(x)=0.
}
\]

Thus, for large \(\lambda\), major contributions to an oscillatory
integral often arise near critical points of the phase.

This is the basis of the stationary phase method.

%%%%%%%%%%%%%%%%%%%%%%%%%%%%%%%%%%%%%%%%%%%%%%%%%%%%%%%%%%%%
\subsection{Restriction Theory Preview}
\label{subsec:restriction-theory-preview}
%%%%%%%%%%%%%%%%%%%%%%%%%%%%%%%%%%%%%%%%%%%%%%%%%%%%%%%%%%%%

A subtle question asks whether the Fourier transform of an \(L^p\)
function can meaningfully be restricted to a lower-dimensional surface
such as

\[
\boxed{
S^{n-1}.
}
\]

Since such a surface has Lebesgue measure zero in \(\R^n\), ordinary
\(L^p\) theory does not automatically provide a restriction.

Fourier restriction theory studies estimates of the form

\[
\boxed{
\|\widehat{f}|_S\|_{L^q(S)}
\leq
C\|f\|_{L^p(\R^n)}.
}
\]

This area connects harmonic analysis to geometry and PDEs.

%%%%%%%%%%%%%%%%%%%%%%%%%%%%%%%%%%%%%%%%%%%%%%%%%%%%%%%%%%%%
\subsection{Wavelets Preview}
\label{subsec:wavelets-preview}
%%%%%%%%%%%%%%%%%%%%%%%%%%%%%%%%%%%%%%%%%%%%%%%%%%%%%%%%%%%%

Fourier analysis localizes well in frequency but complex exponentials are
spread across all of space.

Wavelets provide localization in both position and scale.

A typical family has the form

\[
\boxed{
\psi_{j,k}(x)
=
2^{j/2}\psi(2^jx-k).
}
\]

Here

\[
j
\]

controls scale and

\[
k
\]

controls location.

Wavelets provide another powerful decomposition of function spaces.

%%%%%%%%%%%%%%%%%%%%%%%%%%%%%%%%%%%%%%%%%%%%%%%%%%%%%%%%%%%%
\subsection{Harmonic Analysis on Groups}
\label{subsec:harmonic-analysis-groups}
%%%%%%%%%%%%%%%%%%%%%%%%%%%%%%%%%%%%%%%%%%%%%%%%%%%%%%%%%%%%

Classical Fourier analysis can be generalized from

\[
\R^n
\]

and the circle to locally compact groups.

The group structure provides notions of

\[
\boxed{
\text{translation},
\quad
\text{convolution},
\quad
\text{representation}.
}
\]

On abelian groups, characters generalize complex exponentials.

On nonabelian groups, Fourier analysis becomes closely related to
representation theory.

%%%%%%%%%%%%%%%%%%%%%%%%%%%%%%%%%%%%%%%%%%%%%%%%%%%%%%%%%%%%
\subsection{Common Errors}
\label{subsec:harmonic-analysis-common-errors}
%%%%%%%%%%%%%%%%%%%%%%%%%%%%%%%%%%%%%%%%%%%%%%%%%%%%%%%%%%%%

\subsubsection{Assuming Every \(L^1\) Operator Is Strong Type \((1,1)\)}

Important operators such as the Hardy--Littlewood maximal operator and
many singular integrals may only satisfy weak endpoint estimates.

%%%%%%%%%%%%%%%%%%%%%%%%%%%%%%%%%%%%%%%%%%%%%%%%%%%%%%%%%%%%
\subsubsection{Confusing Weak-Type Estimates with Weak Convergence}

Weak type \((p,q)\) refers to estimates on level sets.

It is unrelated to weak convergence in Banach spaces.

%%%%%%%%%%%%%%%%%%%%%%%%%%%%%%%%%%%%%%%%%%%%%%%%%%%%%%%%%%%%
\subsubsection{Ignoring Cancellation}

A singular kernel may define a bounded operator because of cancellation,
even when its absolute value is not integrable.

%%%%%%%%%%%%%%%%%%%%%%%%%%%%%%%%%%%%%%%%%%%%%%%%%%%%%%%%%%%%
\subsubsection{Treating Principal Value as an Ordinary Improper Integral}

Principal value removes singular regions symmetrically or according to a
specified truncation process.

The cancellation is part of the definition.

%%%%%%%%%%%%%%%%%%%%%%%%%%%%%%%%%%%%%%%%%%%%%%%%%%%%%%%%%%%%
\subsubsection{Assuming \(L^\infty=\operatorname{BMO}\)}

We have

\[
L^\infty
\subsetneq
\operatorname{BMO}.
\]

BMO controls average oscillation, not pointwise magnitude.

%%%%%%%%%%%%%%%%%%%%%%%%%%%%%%%%%%%%%%%%%%%%%%%%%%%%%%%%%%%%
\subsubsection{Forgetting the Endpoint \(p=1\)}

Many operators bounded on

\[
L^p,
\qquad
1<p<\infty,
\]

fail to be bounded on \(L^1\).

Endpoint theory often requires Hardy spaces or weak-type estimates.

%%%%%%%%%%%%%%%%%%%%%%%%%%%%%%%%%%%%%%%%%%%%%%%%%%%%%%%%%%%%
\subsubsection{Assuming \(L^2\) Arguments Automatically Extend to \(L^p\)}

Plancherel makes many Fourier multiplier estimates easy on \(L^2\).

For

\[
p\neq2,
\]

additional smoothness, interpolation, or singular integral theory may be
required.

%%%%%%%%%%%%%%%%%%%%%%%%%%%%%%%%%%%%%%%%%%%%%%%%%%%%%%%%%%%%
\subsubsection{Ignoring Scaling}

An estimate that violates the natural scaling of an operator is usually
impossible.

Scaling should be checked before attempting a long proof.

%%%%%%%%%%%%%%%%%%%%%%%%%%%%%%%%%%%%%%%%%%%%%%%%%%%%%%%%%%%%
\subsubsection{Confusing Harmonic Functions with Harmonic Oscillations}

A harmonic function satisfies

\[
\Delta u=0.
\]

This meaning is related historically and analytically to Fourier
harmonics, but the terms are not identical definitions.

%%%%%%%%%%%%%%%%%%%%%%%%%%%%%%%%%%%%%%%%%%%%%%%%%%%%%%%%%%%%
\subsubsection{Assuming Pointwise Recovery Holds Everywhere}

The Lebesgue Differentiation Theorem gives recovery almost everywhere,
not necessarily at every point.

%%%%%%%%%%%%%%%%%%%%%%%%%%%%%%%%%%%%%%%%%%%%%%%%%%%%%%%%%%%%
\subsection{Applications}
\label{subsec:harmonic-analysis-applications}
%%%%%%%%%%%%%%%%%%%%%%%%%%%%%%%%%%%%%%%%%%%%%%%%%%%%%%%%%%%%

\begin{applicationbox}
\textbf{Partial Differential Equations.}
Maximal functions, singular integrals, Fourier multipliers, and
Littlewood--Paley decompositions provide fundamental estimates for PDEs.

\medskip

\textbf{Potential Theory.}
Harmonic functions, Poisson kernels, Riesz potentials, and singular
integrals describe equilibrium and source fields.

\medskip

\textbf{Complex Analysis.}
Harmonic conjugates, Hardy spaces, boundary values, and Hilbert transforms
connect real-variable and complex-variable methods.

\medskip

\textbf{Signal Processing.}
Frequency localization, filters, wavelets, and singular transforms
separate and analyze features at different scales.

\medskip

\textbf{Image Analysis.}
Multiscale decompositions and singular operators detect edges, textures,
and geometric structures.

\medskip

\textbf{Probability.}
Maximal inequalities, martingale analogues, and Fourier estimates play
major roles in stochastic analysis.

\medskip

\textbf{Geometric Analysis.}
Oscillatory integrals and restriction estimates connect Fourier behavior
with curvature and geometry.

\medskip

\textbf{Number Theory.}
Exponential sums and oscillatory cancellation are central tools in
analytic number theory.

\medskip

\textbf{Fluid Mechanics.}
Riesz transforms and singular integral operators appear in pressure and
velocity formulations of fluid equations.

\medskip

\textbf{Quantum Mechanics.}
Frequency decomposition, uncertainty, oscillatory integrals, and spectral
methods describe wave behavior and propagation.
\end{applicationbox}

%%%%%%%%%%%%%%%%%%%%%%%%%%%%%%%%%%%%%%%%%%%%%%%%%%%%%%%%%%%%
\subsection{Exercises}
\label{subsec:harmonic-analysis-exercises}
%%%%%%%%%%%%%%%%%%%%%%%%%%%%%%%%%%%%%%%%%%%%%%%%%%%%%%%%%%%%

\begin{exercise}[1]
Show that translation preserves the \(L^p\) norm:

\[
\|\tau_yf\|_p=\|f\|_p.
\]
\end{exercise}

\begin{exercise}[1]
Compute the \(L^p\) scaling of

\[
f_t(x)=f(tx)
\]

on \(\R^n\).
\end{exercise}

\begin{exercise}[1]
Show that modulation preserves \(L^p\) norms.
\end{exercise}

\begin{exercise}[1]
Verify that convolution is commutative whenever the relevant integrals
exist.
\end{exercise}

\begin{exercise}[1]
Show that

\[
\int_{\R^n}\phi_\varepsilon(x)\,dx=1
\]

for the standard rescaling of an approximate identity.
\end{exercise}

\begin{exercise}[2]
Prove

\[
\|f*g\|_p
\leq
\|f\|_p\|g\|_1.
\]
\end{exercise}

\begin{exercise}[2]
Let \(f\in C_c(\R^n)\). Show directly that

\[
f*\phi_\varepsilon\to f
\]

uniformly for a standard approximate identity.
\end{exercise}

\begin{exercise}[2]
Explain why mollification smooths a locally integrable function.
\end{exercise}

\begin{exercise}[2]
For

\[
f=\mathbf{1}_{[0,1]},
\]

describe qualitatively the Hardy--Littlewood maximal function \(Mf\).
\end{exercise}

\begin{exercise}[2]
Show that

\[
|A_rf(x)|
\leq
Mf(x).
\]
\end{exercise}

\begin{exercise}[2]
Explain the difference between a strong-type \((1,1)\) estimate and a
weak-type \((1,1)\) estimate.
\end{exercise}

\begin{exercise}[2]
Show that every \(L^\infty\) function belongs to BMO.
\end{exercise}

\begin{exercise}[2]
Show that adding a constant to a function does not change its BMO
seminorm.
\end{exercise}

\begin{exercise}[2]
Verify that harmonic functions satisfy the mean-value property in a
simple radial example.
\end{exercise}

\begin{exercise}[2]
Show that the Poisson kernel satisfies

\[
\frac{1}{2\pi}
\int_{-\pi}^{\pi}
P_r(\theta)\,d\theta=1.
\]
\end{exercise}

\begin{exercise}[3]
Prove the \(L^p\) convergence of approximate identities for

\[
1\leq p<\infty.
\]
\end{exercise}

\begin{exercise}[3]
Study the covering argument used to prove the weak \((1,1)\) estimate for
the Hardy--Littlewood maximal operator.
\end{exercise}

\begin{exercise}[3]
Use the maximal theorem to prove the Lebesgue Differentiation Theorem.
\end{exercise}

\begin{exercise}[3]
Use interpolation between

\[
L^1\to L^\infty
\]

and

\[
L^2\to L^2
\]

to motivate the Hausdorff--Young inequality.
\end{exercise}

\begin{exercise}[3]
Show formally that the Hilbert transform has Fourier multiplier

\[
-i\,\operatorname{sgn}(\xi).
\]
\end{exercise}

\begin{exercise}[3]
Use Plancherel's Theorem to show that the Hilbert transform is bounded on
\(L^2(\R)\).
\end{exercise}

\begin{exercise}[3]
Show formally that

\[
H(Hf)=-f
\]

on an appropriate class of functions.
\end{exercise}

\begin{exercise}[3]
Verify that the Riesz-transform multiplier has absolute value at most
\(1\).
\end{exercise}

\begin{exercise}[3]
Use Plancherel's Theorem to obtain \(L^2\) boundedness of the Riesz
transforms.
\end{exercise}

\begin{exercise}[3]
Check the scaling relation

\[
\frac1q
=
\frac1p-\frac{\alpha}{n}
\]

for the Riesz potential \(I_\alpha\).
\end{exercise}

\begin{exercise}[3]
Show that the Poisson extension of

\[
e^{in\theta}
\]

is

\[
r^{|n|}e^{in\theta}.
\]
\end{exercise}

\begin{exercise}[3]
Use Fourier series to verify that the Poisson extension is harmonic in the
unit disk.
\end{exercise}

\begin{exercise}[4]
Prove the Hardy--Littlewood maximal theorem in one dimension.
\end{exercise}

\begin{exercise}[4]
Study the Marcinkiewicz Interpolation Theorem and use it to derive
\(L^p\) boundedness of the maximal operator from endpoint information.
\end{exercise}

\begin{exercise}[4]
Prove the Riesz--Thorin Interpolation Theorem using complex interpolation.
\end{exercise}

\begin{exercise}[4]
Prove the Hausdorff--Young inequality using interpolation.
\end{exercise}

\begin{exercise}[4]
Study the proof of \(L^p\) boundedness of the Hilbert transform for

\[
1<p<\infty.
\]
\end{exercise}

\begin{exercise}[4]
Study a Calder\'on--Zygmund decomposition of an \(L^1\) function at height
\(\lambda\).
\end{exercise}

\begin{exercise}[4]
Use Calder\'on--Zygmund decomposition to investigate weak endpoint
estimates for singular integrals.
\end{exercise}

\begin{exercise}[4]
Study the proof of the John--Nirenberg inequality for BMO.
\end{exercise}

\begin{exercise}[4]
Investigate the duality

\[
(H^1)^*
\cong
\operatorname{BMO}.
\]
\end{exercise}

\begin{exercise}[4]
Construct a Littlewood--Paley decomposition and prove the \(L^2\) square
function identity up to equivalence constants.
\end{exercise}

\begin{exercise}[4]
Study a proof of the Hardy--Littlewood--Sobolev inequality.
\end{exercise}

\begin{exercise}[4]
Investigate stationary phase for a one-dimensional oscillatory integral
with a nondegenerate critical point.
\end{exercise}

%%%%%%%%%%%%%%%%%%%%%%%%%%%%%%%%%%%%%%%%%%%%%%%%%%%%%%%%%%%%
\subsection{Conceptual Summary}
\label{subsec:harmonic-analysis-summary}
%%%%%%%%%%%%%%%%%%%%%%%%%%%%%%%%%%%%%%%%%%%%%%%%%%%%%%%%%%%%

Harmonic analysis extends Fourier analysis from explicit frequency
decomposition to a broad theory of operators, averages, oscillation, and
function spaces.

Translation is

\[
\boxed{
\tau_yf(x)=f(x-y).
}
\]

Convolution is

\[
\boxed{
(f*g)(x)
=
\int_{\R^n}f(x-y)g(y)\,dy.
}
\]

Approximate identities recover functions through

\[
\boxed{
f*\phi_\varepsilon\to f.
}
\]

The Hardy--Littlewood maximal function is

\[
\boxed{
Mf(x)
=
\sup_{r>0}
\frac1{|B(x,r)|}
\int_{B(x,r)}|f(y)|\,dy.
}
\]

It satisfies

\[
\boxed{
M:L^1\to L^{1,\infty}
}
\]

in the weak sense and

\[
\boxed{
M:L^p\to L^p,
\qquad
1<p\leq\infty.
}
\]

The Lebesgue Differentiation Theorem gives

\[
\boxed{
\lim_{r\to0}
\frac1{|B(x,r)|}
\int_{B(x,r)}f(y)\,dy
=
f(x)
}
\]

for almost every \(x\).

The Hilbert transform is

\[
\boxed{
Hf(x)
=
\frac1\pi
\operatorname{p.v.}
\int_{\R}
\frac{f(y)}{x-y}\,dy,
}
\]

with Fourier multiplier

\[
\boxed{
-i\,\operatorname{sgn}(\xi).
}
\]

Calder\'on--Zygmund theory generalizes such singular integral operators.

Littlewood--Paley theory decomposes functions by frequency scale:

\[
\boxed{
f\sim\sum_j\Delta_jf.
}
\]

Hardy spaces and BMO provide endpoint function spaces adapted to
oscillation and cancellation.

Harmonic functions satisfy

\[
\boxed{
\Delta u=0
}
\]

and possess the mean-value property.

The Poisson kernel produces harmonic extensions and acts as an
approximate identity.

Harmonic analysis therefore studies how

\[
\boxed{
\text{local behavior},
\quad
\text{global norms},
\quad
\text{frequency},
\quad
\text{geometry},
\quad
\text{oscillation}
}
\]

interact.

%%%%%%%%%%%%%%%%%%%%%%%%%%%%%%%%%%%%%%%%%%%%%%%%%%%%%%%%%%%%
\subsection{Section Connections}
\label{subsec:harmonic-analysis-connections}
%%%%%%%%%%%%%%%%%%%%%%%%%%%%%%%%%%%%%%%%%%%%%%%%%%%%%%%%%%%%

\chapterconnection{
Harmonic Analysis builds directly on Fourier Analysis, Measure Theory,
Lebesgue Integration, and Functional Analysis. Maximal functions and
differentiation theory explain almost-everywhere recovery from local
averages. Interpolation transfers operator estimates between function
spaces. Singular integral operators such as the Hilbert and Riesz
transforms provide fundamental tools for Partial Differential Equations
and Potential Theory. Littlewood--Paley decompositions connect frequency
localization to Sobolev spaces and regularity. Hardy spaces and BMO
provide endpoint replacements for ordinary \(L^p\) theory. Harmonic
functions and Poisson kernels connect the subject to Complex Analysis and
Potential Theory. Oscillatory integrals and restriction theory lead
toward Geometric Analysis, while generalized Fourier transforms and
singular kernels connect naturally to Distribution Theory and Operator
Theory.
}

%%%%%%%%%%%%%%%%%%%%%%%%%%%%%%%%%%%%%%%%%%%%%%%%%%%%%%%%%%%%
% SECTION SOURCE TRACKING
%%%%%%%%%%%%%%%%%%%%%%%%%%%%%%%%%%%%%%%%%%%%%%%%%%%%%%%%%%%%

%------------------------------------------------------------
% 5.11 Harmonic Analysis
%------------------------------------------------------------
%
% Source basis:
%   - Standard introductory and beginning graduate harmonic analysis.
%   - Standard real-variable harmonic analysis.
%   - Standard Fourier and singular-integral theory.
%
% Used for:
%   - translation, dilation, and modulation;
%   - convolution;
%   - Young's convolution inequality;
%   - approximate identities;
%   - mollifiers;
%   - averaging operators;
%   - Hardy--Littlewood maximal functions;
%   - weak-type and strong-type estimates;
%   - maximal theorem;
%   - covering lemmas;
%   - Lebesgue differentiation;
%   - interpolation;
%   - Riesz--Thorin interpolation;
%   - Marcinkiewicz interpolation;
%   - Hausdorff--Young inequality;
%   - singular integrals;
%   - principal values;
%   - Hilbert transform;
%   - Calderon--Zygmund theory;
%   - Riesz transforms;
%   - Fourier multipliers;
%   - frequency localization;
%   - Littlewood--Paley theory;
%   - Hardy spaces;
%   - atomic decomposition;
%   - BMO;
%   - John--Nirenberg inequality;
%   - Hardy--BMO duality;
%   - harmonic functions;
%   - mean-value property;
%   - maximum principle;
%   - Poisson kernels;
%   - harmonic extensions;
%   - harmonic conjugates;
%   - subharmonic functions;
%   - fractional integration;
%   - Hardy--Littlewood--Sobolev principle;
%   - scaling;
%   - oscillatory integrals;
%   - stationary phase;
%   - restriction theory preview;
%   - wavelets preview;
%   - harmonic analysis on groups.
%
% Notes:
%   - Fourier Analysis develops Fourier series and Fourier transforms
%     before the operator-theoretic methods used here.
%   - Functional Analysis supplies the Banach-space and Hilbert-space
%     framework for operator estimates.
%   - Measure Theory and Lebesgue Integration provide the L^p spaces
%     and almost-everywhere language used throughout.
%   - Potential Theory later develops harmonic and subharmonic
%     functions in greater depth.
%   - Distribution Theory extends singular kernels and Fourier
%     transforms to generalized functions.
%   - Geometric Analysis and PDE theory use many of the estimates
%     introduced here.
%
%%%%%%%%%%%%%%%%%%%%%%%%%%%%%%%%%%%%%%%%%%%%%%%%%%%%%%%%%%%%